% Evoluzione dal test manuale (docs/VM-test.md) a vm-test.sh e vnic.sh (sezione 4.1).
\begin{table}[H]
    \centering
    \footnotesize
    \caption{Come sono cambiati i principali aspetti dal test manuale a \texttt{vnic.sh}}
    \label{tab:evoluzione}
    \begin{tabular}{@{}>{\raggedright\arraybackslash}p{2.5cm}>{\raggedright\arraybackslash}p{3.7cm}>{\raggedright\arraybackslash}p{3.7cm}>{\raggedright\arraybackslash}p{3.7cm}@{}}
        \toprule
        \textbf{Aspetto} & \textbf{Test manuale} & \textbf{\texttt{vm-test.sh}} & \textbf{\texttt{vnic.sh}} \\
        \midrule
        infrastruttura dell'host & comandi eseguiti a mano & creata dallo script & creata dallo script \\
        \addlinespace
        programma XDP & filtro UDP fisso & filtro UDP fisso & scelto dall'utente (C, file oggetto o P4) \\
        \addlinespace
        configurazione VM & comandi manuali nel terminale della VM & comandi eseguiti da cloud-init al boot & comandi eseguiti da cloud-init al boot \\
        \addlinespace
        immagine della VM & immagine base modificata & overlay dell'immagine base & overlay dell'immagine base, scaricata se assente \\
        \addlinespace
        attesa del boot & login manuale & lettura automatica del log della VM & lettura automatica del log della VM \\
        \addlinespace
        verifica & \texttt{tcpdump}, \texttt{ping} e \texttt{nc} a mano & comando \texttt{run}, con esito PASS o FAIL & separata dal tool: \texttt{traffic.sh} e \texttt{verify.sh} degli esempi \\
        \bottomrule
    \end{tabular}
\end{table}
